% ============ §33 Some other applications (原书 pp.261-265) ============
\section{Some other applications}\label{sec:33}\index{dimension of intersection, section 33}\index{integral closure, finiteness of, section 33}\index{Nagata ring, section 33}\index{Mori-Nagata integral closure theorem, section 33}

\medskip
\noindent\textbf{Dimension of intersection}

Let $k$ be a field and $V$ and $W$ irreducible algebraic varieties in affine
$n$-space over $k$. (Here one may either assume that $k$ is algebraically closed
and identify the varieties with the corresponding subsets of $k^n$, or take the
scheme-theoretic viewpoint.) Then it is well known that every irreducible
component of $V \cap W$ has dimension $\ge \dim V + \dim W - n$. Algebraically, this
is equivalent to the following theorem.

Let $P$ and $P'$ be two prime ideals in the polynomial ring $R = k[X_1, \ldots, X_n]$
over a field $k$, and let $Q$ be a minimal prime divisor of $P + P'$. Then
\[
  \dim R/Q \ge \dim(R/P) + \dim(R/P') - n,
\]
or equivalently,
\[
  (*) \quad \htop Q \le \htop P + \htop P'.
\]
The idea of the proof consists of transforming the intersection $V \cap W$ in $k^n$
into the intersection $\Delta \cap (V \times W)$ in $k^{2n}$, where $\Delta$ is the diagonal, and
availing oneself of the fact that $\Delta$ is defined by $n$ equations
$x_i - y_i = 0$, for $i = 1, \ldots, n$ (see [M], p.~93).

Now in algebraic geometry, the theorem remains true if one replaces affine
$n$-space by a non-singular (smooth) variety; namely, if $V$ and $W$ are
irreducible subvarieties of a non-singular variety $U$, then every irreducible
component of $V \cap W$ has dimension $\ge \dim V + \dim W - \dim U$. Algebraically,
the inequality (*) still holds if $R$ is an arbitrary regular local ring containing
a field. One can easily reduce to the case where $R$ is complete, and then by
I.~S.~Cohen's structure theorem $R$ is isomorphic to $k[[X_1, \ldots, X_n]]$, and so one
can apply the same diagonal trick. Thus it is clear that (*) holds in an
arbitrary regular local ring of equal characteristic. How about the unequal
characteristic case? If $R$ is unramified, then its completion $R^*$ is a formal
power series ring over a DVR by Theorem~\hyperref[thm:29.8]{29.8}, and a slight modification of the
argument used in the case of $k[X_1, \ldots, X_n]$ suffices. When $R$ is ramified, by
Theorem~\hyperref[thm:29.9]{29.9}, $R^*$ is of the form $D[X_1, \ldots, X_n]/(f)$, where $D$ is a complete
DVR and $f$ is an Eisenstein polynomial. Using this, and applying his deep
results on intersection multiplicity, J.-P.~Serre proved the inequality (*) for
general regular local rings $R$ in Chapter V of his book [Se]. We recommend this
excellent book to the reader.

\medskip
\noindent\textbf{Integral closure of a Noetherian integral domain}

Let $A$ be a Noetherian integral domain with field of fractions $K$, and let $A'$
denote the integral closure of $A$ in $K$ (the so-called \emph{derived normal ring} of
$A$). Is $A'$ a finite module over $A$? This is a difficult question, and the answer
is no in general. When $A$ is finitely generated over a field $k$ (the case
encountered in algebraic geometry) it is easy to prove finiteness. We need the
following two lemmas.

\begin{lemma}\label{lem:33.1}
Let $A$ be a Noetherian normal integral domain with field of fractions $K$;
suppose that $L$ is a finite separable extension of $K$, and let $A'$ be the
integral closure of $A$ in $L$. Then $A'$ is a finite $A$-module.
\end{lemma}

\begin{proof}
By enlarging $L$ if necessary we can assume that $L$ is a Galois extension of $K$.
Write $G = \{\sigma_i \mid 1 \le i \le n\}$ for the Galois group of $L/K$, where
$n = [L : K]$, and let $y_1, \ldots, y_n$ be elements of $A'$ which form a basis of $L$
over $K$. If $z \in A'$ and $z = \sum_j c_jy_j$ with $c_j \in K$, then
$\sigma_iz = \sum_j c_j\sigma_iy_j$ for $i = 1, \ldots, n$, and hence
$c_j = C_j/D$, where $D = \det(\sigma_iy_j)$ and $C_j \in A'$. Putting $d = D^2$ we see
$d \in K$. In fact it is easy to see that $d = \det(\tr_{L/K}(y_iy_j))$ is the
discriminant of the separable $K$-algebra $L$ (see p.~198). It follows that
$d \ne 0$ and $dc_j \in A' \cap K = A$ for all $j$. Therefore $A'$ is contained in the
finite $A$-module $\sum_j Ad^{-1}y_j$, so that $A'$ itself is finite over $A$.
\end{proof}

\begin{lemma}[\textbf{Normalisation theorem of E.~Noether}]\label{lem:33.2}
Let $A = k[x_1, \ldots, x_n]$ be a finitely generated algebra over a field $k$; then
there exist $y_1, \ldots, y_r \in A$ which are algebraically independent over $k$ such
that $A$ is integral over $k[y_1, \ldots, y_r]$.
\end{lemma}

\begin{proof}
Here we assume that $k$ is an infinite field, referring the reader to [M], (14.G)
or [N1], (14.4) for the general case. Suppose $x_1, \ldots, x_n$ are algebraically
dependent over $k$, and let $f(x_1, \ldots, x_n) = 0$ be a relation, where
$f(X_1, \ldots, X_n)$ is a non-zero polynomial with coefficients in $k$. Write $d$ for
the degree of $f$ and let $f_d(X_1, \ldots, X_n)$ be the homogeneous part of $f$ of
degree $d$. Take $c_1, \ldots, c_{n-1} \in k$ such that
$f_d(c_1, \ldots, c_{n-1}, 1) \ne 0$, and set $y_i = x_i - c_ix_n$ for $i = 1, \ldots, n - 1$. Then
\[
  0 = f(x_1, \ldots, x_n) = f(y_1 + c_1x_n, \ldots, y_{n-1} + c_{n-1}x_n, x_n)
      = f_d(c_1, \ldots, c_{n-1}, 1)x_n^d + g_1x_n^{d-1} + \cdots + g_d,
\]
with $g_i \in k[y_1, \ldots, y_{n-1}]$, so that $x_n$ is integral over
$k[y_1, \ldots, y_{n-1}]$. Then $x_i = y_i + c_ix_n$, for $i = 1, \ldots, n - 1$, are also
integral over $k[y_1, \ldots, y_{n-1}]$, hence $A$ is integral over
$k[y_1, \ldots, y_{n-1}]$. Thus the assertion is proved by induction on $n$.
\end{proof}

Now let $A$ be a finitely generated integral domain over $k$, with field of
fractions $K$ and derived normal ring $A'$. Take $y_1, \ldots, y_r \in A$ as in
Lemma~\hyperref[lem:33.2]{2}, so that $A'$ is also the integral closure of $k[y_1, \ldots, y_r]$ in $K$.
Set $K' = k(y_1, \ldots, y_r)$. Then $K$ is a finite algebraic extension of $K'$. If this
extension is separable then $A'$ is finite over $k[y_1, \ldots, y_r]$ by
Lemma~\hyperref[lem:33.1]{1}, hence is also finite over $A$. If $K$ is inseparable over $K'$, let
$p = \chark K$. Then there is a purely inseparable finite extension $K''$ of $K'$ such
that $K(K'')$ is separable over $K''$. Therefore it suffices to prove that the
integral closure of $k[y_1, \ldots, y_r]$ in $K''$ is finite over it. But $K''$ is
contained in a field $L$ which is obtained by adjoining to $K'$ the $q$-th roots of
a finite number of elements $a_1, \ldots, a_s$ of $k$ and also the $q$-th roots of
$y_1, \ldots, y_r$, where $q$ is a sufficiently high power of $p$. Then the integral
closure of $k[y_1, \ldots, y_r]$ in $L$ is $k'[y_1^{1/q}, \ldots, y_r^{1/q}]$, where
$k' = k(a_1^{1/q}, \ldots, a_s^{1/q})$, and it is clear that $k'[y_1^{1/q}, \ldots, y_r^{1/q}]$ is
finite over $k[y_1, \ldots, y_r]$. This completes the proof of finiteness of $A'$ over
$A$.

When $A$ is a complete Noetherian local domain one can prove the finiteness of
$A'$ along the same line as above. Using Theorem~\hyperref[thm:29.5]{29.5}, (iii), instead of
Lemma~\hyperref[lem:33.2]{2}, one reduces to proving the finiteness of the integral closure of a
complete regular local ring $A$ in a finite extension $L$ of the field of fractions
$K$ of $A$. If $\chark K = 0$ this is proved by Lemma~\hyperref[lem:33.1]{1}, so that we can assume
$\chark K = p > 0$. Then $A = k[[X_1, \ldots, X_n]]$ is a formal power series ring over a
field $k$. We can also assume, as in the above proof, that $L$ is purely
inseparable over $K$, so that there is a power $q = p^m$ of $p$ such that
$L \subset K^{1/q}$. But there is one problem. Since a formal power series has
infinitely many coefficients, it may not be possible to find a finite extension
$k_1$ of $k$ such that $L \subset k_1((Y))$, where $Y = (Y_1, \ldots, Y_n)$,
$Y_i = X_i^{1/q}$ and $k_1((Y))$ denotes the field of fractions of $k_1[[Y]]$. One can
overcome this difficulty either by an argument of Nagata in [N1, p.~113], or
(following J.~Tate) by induction on $n$ as follows:

We may assume that $x = X_1^{1/q} \in L$ ($1 \le i \le n$). Since $A$ is normal we have
$A' = \{f \in L \mid f^q \in A\}$. Set $P = X_1A$, $Q = Y_1A'$; then
$Q = (\xi \in L \mid \xi^q \in P)$, so that $Q$ is the only prime ideal lying over $P$. Now
$A_P$ and $A'_Q$ are DVRs by Theorem~\hyperref[thm:11.2]{11.2}, (3), and their fields of fractions
are $L$ and $K$ respectively. Let $K'$ and $K_0$ be their residue fields; then
$[K' : K_0] \le [L : K]$ by Ex.~\hyperref[ex:10.8-ex]{10.8}. Since $A'/Q$ is contained in the integral
closure of $A/P = k[[X_2, \ldots, X_n]]$ in $K'$, by the induction hypothesis, $A'/Q$ is
finite over $A/P$. Since
\[
  Q^i/Q^{i+1} = Y_1^iA'/Y_1^{i+1}A' \cong A'/Q \quad \text{and} \quad Q^q = PA'
\]
we see that $A'/PA'$ is finite over $A/P$. Moreover, $A'$ is separated in the $P$-adic
topology (which is the same as the $Y_1$-adic topology), because
$A' \subset k^{1/q}[[Y_1, \ldots, Y_n]]$. Since $A$ is $P$-adically complete, $A'$ is finite over
$A$ by Theorem~\hyperref[thm:8.4]{8.4}.

From this result it is easy to derive the following theorem: \emph{If $A$ is a Noetherian
local ring whose completion $A^*$ is reduced, then the integral closure $A'$ of $A$ in its
total ring of fractions is finite over $A$}. See [M] p.~237.

On the other hand, if $A$ is not reduced and if the maximal ideal $m$ contains a
regular element (that is, a non-zero-divisor), then $A'$ is not finite over $A$
(Krull [2]). In fact, if $x \ne 0$ is nilpotent, take a regular element $s$ such that
$x \notin sA$; (it is possible to find such an $s$ since $\bigcap_i m^i = (0)$). Then the
elements $x/s^j$ (for $j = 1, 2, 3, \ldots$) belong to $A'$. If $A'$ is finite over $A$ then
there must be some integer $n$ such that $s^n(x/s^j) \in A$ for all $j$. But then
$x \in \bigcap_{j \ge 0}s^jA = (0)$, a contradiction.

Suppose $(A, m)$ is a one-dimensional Noetherian local integral domain. Then $A^*$
is reduced if and only if $A'$ is finite over $A$ (Krull [2]). In fact, if $A'$ is finite
over $A$ then it is a semilocal ring, and for each maximal ideal $P$ of $A'$ the
local ring $(A')_P$ is a DVR. The $\rad(A')$-adic topology of $A'$ coincides with the
$m$-adic topology, and the completion $A'^*$ of $A'$ with respect to this topology is
a direct product of complete DVRs by Theorem~\hyperref[thm:8.15]{8.15}, hence is reduced. On
the other hand it coincides with $A^* \otimes_A A'$ by Theorem~\hyperref[thm:8.7]{8.7}. Since
$0 \to A \to A'$ is exact, $0 \to A^* \to A^* \otimes_A A'$ is also exact. Therefore $A^*$ is
reduced. The converse holds, as already mentioned, without the restriction on
dimension.

Rees [9] proved that a reduced Noetherian local ring $A$ has reduced completion
$A^*$ if and only if for every finite subset $T$ of the total ring of fractions $K$
of $A$, the integral closure of $A[t]$ in $K$ is finite over $A[t]$.

Akizuki [1] constructed the first example of a one-dimensional Noetherian local
integral domain with non-reduced completion (see also Larfeldt-Lech [1]). To
avoid such pathology, Nagata [N1] defined and studied the class of
pseudo-geometric rings, which were called `anneaux universellement japonais'
by Grothendieck ([G1], [G2]). These are now known as `Nagata rings' ([M],
[B9]). A Noetherian ring $A$ is called a \emph{Nagata ring}\index{Nagata ring} if for every prime
ideal $P$ of $A$ and for every finite extension field $L$ of the field of fractions
$\kappa(P)$ of $A/P$, the integral closure of $A/P$ in $L$ is finite over $A/P$. For the
basic properties of Nagata rings, see [M], \S31. An alternative definition is the
following: a Noetherian ring $A$ is a Nagata ring if (1) for every maximal ideal
$m$ the formal fibres of $A_m$ are geometrically reduced, and (2) for every finite
$A$-algebra $B$ which is an integral domain, the set $\Nor(B) = \{P \in \Spec B \mid B_P
\text{ is normal}\}$ is open in $\Spec B$. The equivalence of these two definitions
can easily be proved from [G2], (7.6.4) and (7.7.2).

Although the integral closure $A'$ of a Noetherian integral domain $A$ may fail to
be finite over $A$, it is a Krull ring by the theorem of Mori-Nagata mentioned in
\S\ref{sec:12}. Because of its importance we quote here the theorem in full.

\begin{theorem}[\textbf{Mori-Nagata integral closure theorem}]\label{thm:33.1}
Let $A$ be a Noetherian integral domain and let $A'$ be its derived normal ring.
Then (1) $A'$ is a Krull ring, and (2) for every prime ideal $P$ of $A$ there are
only finitely many prime ideals $P'$ of $A'$ lying over $P$, and for each such $P'$
the field of fractions $\kappa(P')$ of $A'/P'$ is finite over $\kappa(P)$.
\end{theorem}

For a proof, see [N1], (33.10). This proof depends on I.~S.~Cohen's structure
theorem. There are also more recent proofs which do not use the structure
theorem (Nishimura [2], Querr\'e [1], Kiyek [1]).

Note that $A'$ is Noetherian if $\dim A \le 2$. This follows easily from the above
theorem, Theorem~\hyperref[thm:11.7]{11.7} (Krull-Akizuki) and Theorem~\hyperref[thm:12.7]{12.7} (Mori-Nishimura).
When $\dim A = 3$, Nagata constructed a counter-example ([N1], p.~207).

Theorem~\hyperref[thm:28.9]{28.9}, which is due to Grothendieck and is not proved in this book,
was given a new proof by Radu [5]. This interesting proof depends heavily on
I.~S.~Cohen's structure theorem. The same remark applies also to Andr\'e's proof
[1] of the theorem mentioned at the end of \S\ref{sec:32}.

I.~S.~Cohen's structure theorem is also at the basis of the theories of canonical
modules ([HK]) and of dualising complexes (Sharp [3], [5]). Here, the fact that
a complete Noetherian local ring is a quotient of a Gorenstein ring is
important.
